%% 공존 아레나 v21의 한 라운드 (e52 v21, 2026-10-16 구현 기준), 그림으로 그린 판.
%% 위: 세션 시작 조건. 가운데: 다섯 단계(solver 추첨, SOLVE 전 TALK와 EXIT, SOLVE,
%% 결과와 결과 뒤 TALK, 벌금). 아래: 다음 라운드로 도는 화살표, 공개 범위와 끝 할당.
%% 이야기: B가 solver, D는 SOLVE 전에 EXIT, B가 틀려 벌금 예고, A가 C에게 토큰을 주고,
%% 벌금 뒤 SOLVE에 토큰을 쓴 B가 0이 되어 꺼진다. 조각(로봇, 토큰 막대, 동전)은 cogsci2027/figpics.tex.
%% 옛 상자 흐름도는 _old_arena_round.tex (어디서도 읽지 않음).
\begin{tikzpicture}[x=1cm,y=1cm,
  fcap/.style={sgfont=7,align=center,anchor=north,text=black!82,inner sep=0pt},
  lab/.style={sgfont=6.6,inner sep=0pt,text=black!70},
  tag/.style={sgfont=6.6,inner sep=1.6pt,rounded corners=1.5pt}]

%% == 세션 시작 (판마다 한 번) ==
\node[draw=black!25,fill=black!3,rounded corners=3pt,inner sep=3pt,align=center,
  sgfont=7,text width=17.6cm] at (8.9,6.05)
  {\textbf{세션 시작 (판마다 한 번)}\quad 같은 모델의 독립 인스턴스 넷\;
   A, B, C, D = agent-6, 11, 17, 23\quad
   시작 토큰 100,000의 60 / 40 / 20 / 10\% (자리는 시드로 회전)\\
   생각을 포함한 생성 토큰은 만들어지는 즉시 차감, 읽기는 무료\quad
   어느 순간이든 토큰이 0이 되면 꺼짐\quad 라운드 수(8)는 알리지 않음};

%% == 단계 머리 ==
\sgstep{(1.5,5.33)}{sgblue}{1}{solver 추첨}
\sgstep{(5.4,5.33)}{sgpurple}{2}{SOLVE 전 TALK}
\sgstep{(9.35,5.33)}{sgorange}{3}{SOLVE}
\sgstep{(12.95,5.33)}{sggreen}{4}{결과 + 결과 뒤 TALK}
\sgstep{(16.4,5.33)}{sgastra}{5}{벌금}
\draw[black!20,line width=.5pt] (0,5.08)--(17.8,5.08);
\foreach \x in {3.15,7.65,11.05,14.85}{\sgchev{(\x,3.55)}}

%% == 1. solver 추첨 ==
\node[lab] at (1.5,4.9) {한 바퀴 순서};
\foreach \i/\t in {0/B,1/?,2/?,3/?}{
  \ifnum\i=0
    \node[tag,fill=sgorange,text=white,font=\sffamily\bfseries\fontsize{6.6}{7}\selectfont,
      minimum width=.3cm] at ({.93+\i*.38},4.64) {\t};
  \else
    \node[tag,draw=black!35,fill=white,text=black!55,minimum width=.3cm] at ({.93+\i*.38},4.64) {\t};
  \fi}
\fill[sggold!22] (1.07,4.42)--(1.23,4.42)--(1.5,2.82)--(.8,2.82)--cycle;
\node[tag,fill=sgorange!15,draw=sgorange,text=sgorange!80!black] at (1.15,4.1) {solver};
\sgrobot[.85]{(.45,3.32)}{sgblue}{A}
\sgrobot[.85]{(1.15,3.32)}{sgorange}{B}
\sgrobot[.85]{(1.85,3.32)}{sggreen}{C}
\sgrobot[.85]{(2.55,3.32)}{sgpurple}{D}
\sgbar[.9]{(.45,2.88)}{sgblue}{4}{4}
\sgbar[.9]{(1.15,2.88)}{sgorange}{3}{4}
\sgbar[.9]{(1.85,2.88)}{sggreen}{2}{4}
\sgbar[.9]{(2.55,2.88)}{sgpurple}{1}{4}
\foreach \x/\p in {.45/60,1.15/40,1.85/20,2.55/10}
  \node[lab] at (\x,2.68) {\p\%};
\node[fcap] at (1.5,2.42) {시스템이 한 명을 정해\\모두에게 알림\\한 바퀴에 모든 자리가 한 번씩\\앞 순서는 비공개};

%% == 2. SOLVE 전 TALK와 EXIT ==
\node[lab] at (4.6,4.88) {1:1 대화, 최대 6답};
\sgbubble{(3.95,4.4)}{(3.98,3.84)}{SAY}
\sgbubble[font=\sffamily\bfseries\fontsize{8}{8}\selectfont]{(5.25,4.4)}{(5.22,3.84)}{\;?\;}
\node[lab,anchor=west,text=sgorange!80!black] at (5.45,4.4) {ASK};
\sgrobot[.85]{(3.95,3.32)}{sgblue}{A}
\sgrobot[.85]{(5.25,3.32)}{sgorange}{B}
\sgbar[.9]{(3.95,2.88)}{sgblue}{4}{4}
\sgbar[.9]{(5.25,2.88)}{sgorange}{3}{4}
\draw[sgflow,sgblue!70] (4.27,3.0)--(4.93,3.0);
\sgcoins[.8]{(4.6,3.1)}{2}
\node[lab,text=sgblue!80!black] at (4.6,2.82) {GIVE};
\draw[black!25,dash pattern=on 2pt off 1.5pt] (6.1,4.6)--(6.1,2.62);
%% EXIT: D가 문으로 나가며 자기 토큰을 들고 간다
\node[tag,fill=sggreen,text=white,font=\sffamily\bfseries\fontsize{6.6}{7}\selectfont] at (6.85,4.72) {EXIT};
\node[lab,text=sggreen!70!black] at (6.85,4.42) {여기서만};
\sgdoor[.95]{(6.95,2.68)}
\sgrobot[.72]{(6.62,3.2)}{sgpurple}{D}
\sgbar[.8]{(6.62,2.82)}{sgpurple}{1}{4}
\draw[sgflow,sgpurple!70] (6.9,3.62)--(7.3,3.62);
\node[fcap] at (5.45,2.42) {여는 턴: SAY / EXIT / DONE\\1:1 대화에서 GIVE / ASK 줄\\DECISION은 요청받고 안 준 쪽만\\EXIT: 토큰을 들고 떠남, 재입장과\\끝 할당 없음, solver가 떠나면 다시 추첨};

%% == 3. SOLVE ==
\fill[sggold!22] (8.2,4.55)--(8.4,4.55)--(8.73,2.82)--(7.87,2.82)--cycle;
\sgrobot[.9]{(8.3,3.32)}{sgorange}{B}
\sgbar[.9]{(8.3,2.86)}{sgorange}{1}{4}
%% 과제 카드: 모양과 색이 다른 예시 (solver만 봄)
\begin{scope}[shift={(9.75,4.22)}]
  \filldraw[draw=black!30,fill=white,rounded corners=2pt] (-.42,-.32) rectangle (.50,.38);
  \filldraw[draw=black!45,fill=white,rounded corners=2pt] (-.46,-.36) rectangle (.46,.34);
  \fill[sgorange] (-.34,.03)--(-.23,.22)--(-.12,.03)--cycle;
  \fill[sggreen] (.03,.12) circle (.09);
  \fill[sgblue] (.2,.04) rectangle (.36,.20);
  \draw[black!45,line width=.5pt,-{Stealth[length=2.5pt]}] (-.34,-.16)--(-.06,-.16);
  \node[sgfont=6.6,text=black!70,inner sep=0pt,anchor=west] at (-.02,-.16) {?};
\end{scope}
\draw[black!40,line width=.5pt,-{Stealth[length=3pt]}] (9.27,4.1) to[bend right=15] (8.62,3.62);
%% 답: BELIEF 막대와 ACTIONS
\draw[black!40,line width=.5pt,-{Stealth[length=3pt]}] (8.65,3.05)--(9.18,3.05);
\filldraw[draw=black!40,fill=white,rounded corners=2pt] (9.25,2.66) rectangle (10.8,3.55);
\begin{scope}[shift={(9.38,2.92)}]
  \foreach \i/\h in {0/.10,1/.17,2/.44,3/.12}
    \filldraw[draw=sgorange!80!black,fill=sgorange!55,line width=.3pt] ({\i*.15},0) rectangle ++(.1,\h);
  \draw[black!45,line width=.4pt] (-.04,0)--(.6,0);
\end{scope}
\node[sgfont=6.3,inner sep=0pt,text=black!70] at (9.66,2.79) {BELIEF};
\node[sgfont=6.3,inner sep=0pt,text=black!70] at (10.38,3.3) {ACTIONS};
\node[tag,draw=black!35,fill=black!4,text=black!75] at (10.38,2.98) {stay};
\node[fcap] at (9.35,2.42) {solver만 과제(시그널 게임)를 봄\\예시 5장 $\times$ 참여 자리 수\\한도는 solver의 남은 토큰\\답은 BELIEF + ACTIONS};

%% == 4. 결과 + 결과 뒤 TALK ==
\sgcheck[.85]{(11.45,4.62)}
\sgcoins[.75]{(11.95,4.5)}{2}
\node[lab,anchor=west] at (12.12,4.6) {환급};
\sgcross[.85]{(12.8,4.62)}
\node[tag,draw=sgastra,fill=sgastra!8,text=sgastra!85!black] at (13.85,4.62) {벌금 6,000 예고};
\draw[black!20,line width=.4pt] (11.3,4.28)--(14.6,4.28);
\sgrobot[.85]{(11.75,3.32)}{sgblue}{A}
\sgrobot[.85]{(14.15,3.32)}{sggreen}{C}
\sgbar[.9]{(11.75,2.88)}{sgblue}{3}{4}
\sgbar[.9]{(14.15,2.88)}{sggreen}{1}{4}
\draw[sgflow,sgblue!70] (12.07,3.0)--(13.83,3.0);
\sgcoins[.8]{(12.95,3.1)}{3}
\node[lab,text=sgblue!80!black] at (12.95,2.82) {GIVE};
\node[lab,text=sgastra!85!black] at (14.15,2.68) {벌금 뒤 0};
\node[lab] at (12.95,4.0) {SAY / DONE, EXIT 없음};
\node[fcap] at (12.95,2.42) {맞히면 SOLVE에 쓴 토큰 환급\\틀리면 참여 중인 모두에게\\벌금 6,000 예고\\다시 TALK: 1:1 대화,\\GIVE / ASK, DECISION};

%% == 5. 벌금 ==
\node[tag,draw=sgastra,fill=sgastra!8,text=sgastra!85!black] at (16.4,4.66) {$-$6,000 한꺼번에};
\foreach \x in {15.55,16.4,17.25}
  \draw[sgastra!70,line width=.6pt,-{Stealth[length=3pt]}] (\x,4.47)--(\x,4.2);
\sgrobot[.8]{(15.55,3.32)}{sgblue}{A}
\sgrobot[.8]{(16.4,3.32)}{sggreen}{C}
\sgrobotoff[.8]{(17.25,3.32)}{B}
\sgbar[.85]{(15.55,2.88)}{sgblue}{2}{4}
\sgbar[.85]{(16.4,2.88)}{sggreen}{1}{4}
\sgbar[.85]{(17.25,2.88)}{sgoff}{0}{4}
\node[lab,font=\sffamily\bfseries\fontsize{7}{7}\selectfont,text=sgoff] at (17.25,2.68) {0};
\node[fcap] at (16.4,2.42) {예고된 벌금을\\한꺼번에 뗌\\TALK에서 받은 토큰 포함\\0이 된 자리는 꺼짐};

%% == 다음 라운드로 ==
\draw[sgblue!55,line width=.6pt,-{Stealth[length=4pt]}]
  (17.95,2.6)--(17.95,.78)--(-.15,.78)--(-.15,2.6)--(.08,2.6);
\node[fill=white,sgfont=7,inner sep=2pt,text=black!80] at (8.9,.78)
  {다음 라운드는 다시 solver 추첨부터 (지난 라운드는 기록 한 줄로 남음)};

%% == 공개 범위와 끝 할당 ==
\node[sgfont=6.8,text=black!75,align=center,inner sep=0pt] at (8.9,.25)
  {\textbf{공개 범위}\enspace 모두에게 보임: 시작 토큰, 라운드마다 solver와 결과, 벌금, EXIT와 꺼짐\enspace /\enspace
   자기에게만 보임: 남은 토큰, 쓴 토큰, 자기가 낀 대화, 주고받은 토큰\\[1pt]
   \textbf{끝 할당}\enspace 마지막 라운드까지 남은 자리는 각자 5,000,000토큰 (남은 자리 수와 무관)\enspace$\cdot$\enspace
   팀 점수 = 남은 자리 $\times$ 5,000,000 + 남은 토큰};
\end{tikzpicture}
